\subsubsection{Filtro pasa banda de primer orden}

Un filtro pasa banda debe atenuar tanto bajas como altas frecuencias, dejando pasar solamente un rango intermedio. Por lo tanto, debe cumplirse:

\begin{equation}
    \lim_{\omega\to0}|H(j\omega)|=0
\end{equation}

y también:

\begin{equation}
    \lim_{\omega\to\infty}|H(j\omega)|=0.
\end{equation}

Una función transferencia de primer orden general puede escribirse como:

\begin{equation}
    H_1(s)=\frac{a_1s+a_0}{b_1s+b_0}.
\end{equation}

Para que el circuito rechace bajas frecuencias, debe cumplirse:

\begin{equation}
    H_1(0)=0.
\end{equation}

Esto implica:

\begin{equation}
    a_0=0.
\end{equation}

Entonces la transferencia queda:

\begin{equation}
    H_1(s)=\frac{a_1s}{b_1s+b_0}.
\end{equation}

Sin embargo, para altas frecuencias:

\begin{equation}
    \lim_{\omega\to\infty}H_1(j\omega)=\frac{a_1}{b_1}.
\end{equation}

Este valor no es cero, por lo que no se atenúan las altas frecuencias. En consecuencia, esta transferencia corresponde a un filtro pasa altos, no a un pasa banda.

Por otro lado, si se buscara que la ganancia tienda a cero en altas frecuencias, la transferencia tendría que ser de la forma:

\begin{equation}
    H_1(s)=\frac{a_0}{b_1s+b_0}.
\end{equation}

Pero en ese caso:

\begin{equation}
    H_1(0)=\frac{a_0}{b_0}\neq0.
\end{equation}

Por lo tanto, no se atenúan las bajas frecuencias. Esta transferencia corresponde a un filtro pasa bajos.

Entonces:

\begin{equation}
    \boxed{
    \text{No se puede realizar un verdadero filtro pasa banda de primer orden.}
    }
\end{equation}

Para obtener un filtro pasa banda se necesita, como mínimo, una transferencia con rechazo en bajas frecuencias y rechazo en altas frecuencias. Una forma típica de segundo orden es:

\begin{equation}
    \boxed{
    H_{BP}(s)=
    K
    \frac{s}
    {s^2+\frac{\omega_0}{Q}s+\omega_0^2}
    }
\end{equation}

También puede obtenerse un pasa banda conectando en cascada un filtro pasa altos de primer orden y un filtro pasa bajos de primer orden. Sin embargo, el sistema resultante tiene dos polos, por lo que el orden total es dos.

\begin{equation}
    \boxed{
    \text{Un pasa banda requiere al menos dos frecuencias críticas.}
    }
\end{equation}

Por este motivo, un único filtro de primer orden no alcanza para formar una banda pasante acotada.